\clearpage
\phantomsection% 
\addcontentsline{toc}{chapter}{KATA PENGANTAR}
%\thispagestyle{fancy}

\begin{justifying}
	\large \bfseries \centering \MakeUppercase{Kata Pengantar}\linebreak
	
	\normalsize \normalfont \justifying
	Puji syukur ke hadirat Tuhan Yang Maha Esa atas segala rahmat dan karunia-Nya, sehingga penulis dapat menyelesaikan tugas akhir ini dengan baik.
    
    Penyusunan tugas akhir ini tidak terlepas dari bimbingan dan dukungan berbagai pihak. Oleh karena itu, penulis mengucapkan terima kasih yang sebesar-besarnya kepada:
    \begin{enumerate}
        \item Bapak Martin Clinton Tosima Manullang, Ph.D., selaku Dosen Pembimbing yang telah meluangkan waktu dan tenaga untuk mengarahkan penulis hingga tugas akhir ini selesai.
        \item Bapak I Wayan Wiprayoga Wisesa, S.Kom., M.Kom., Radhinka Bagaskara, S.Si.Kom., M.Si., M.Sc., dan Prof. Sarwono Sutikno, Dr, Eng., selaku Dosen Penguji yang telah memberikan masukan, kritik, serta saran konstruktifnya.
        \item Ibu tercinta yang senantiasa memberikan doa yang tak terputus, kasih sayang yang tulus, serta motivasi yang tiada henti kepada penulis dalam setiap langkah.
        \item Kedua kakak dan kedua abang ipar atas dorongan semangat serta dukungannya selama proses penulisan.
        \item Bapak Cahya Wirawan atas penyediaan dan pemeliharaan dataset LibriVox Indonesia di HuggingFace yang sangat membantu penelitian ini.
        \item Seluruh pihak yang tidak dapat disebutkan satu per satu atas segala bantuan dan dukungannya.
    \end{enumerate} 
    
    Penulis menyadari bahwa tugas akhir ini masih jauh dari sempurna. Semoga karya ini dapat memberikan manfaat dan wawasan yang berguna bagi para pembaca.
	
\end{justifying}
\clearpage